\section{Conclusões}%
\label{sec:conclusoes}

O estudo desenvolvido ao longo deste trabalho permitiu explorar de maneira abrangente o problema da Árvore Geradora Mínima (AGM), desde seus fundamentos teóricos até a formulação matemática e obtenção de resultados ótimos por meio de Programação Linear Inteira. A AGM constitui um dos problemas fundamentais em teoria dos grafos e otimização combinatória, e sua resolução exata por métodos matemáticos reforça não apenas sua relevância prática, mas também o valor das técnicas de modelagem aplicadas.

Inicialmente, discutiu-se a estrutura do problema e seus requisitos essenciais: conectividade entre todos os vértices, ausência de ciclos e minimização do custo total associado às arestas selecionadas. Em seguida, abordou-se a formulação matemática clássica, destacando suas limitações, especialmente no que diz respeito às restrições de corte, que possuem caráter exponencial e podem tornar o modelo computacionalmente inviável para instâncias maiores.

A partir disso, introduziu-se a formulação baseada em fluxo (\textit{single-commodity flow}), que representa uma alternativa eficiente e elegante para modelar a conectividade do grafo. Essa abordagem evita explicitamente a necessidade de restrições de corte ao empregar um artifício de conservação de fluxo que garante, de maneira natural, que todos os vértices sejam alcançáveis a partir de um nó raiz. A utilização simultânea de variáveis binárias \(x_{uv}\) e variáveis de fluxo \(f_{uv}\) permitiu expressar de maneira rigorosa a dependência entre a seleção de uma aresta e a possibilidade de transporte de fluxo através dela. As restrições de capacidade desempenharam papel crucial nesse vínculo, assegurando que nenhum fluxo circulasse por arestas inativas.

Os resultados computacionais reforçaram a coerência e a funcionalidade da modelagem. As arestas selecionadas compõem um conjunto de exatamente \(n - 1\) elementos, formando uma estrutura conexa e acíclica, conforme exigido para uma AGM.\@ O custo total obtido, \(C_{\text{total}} = 433\), confirma que o modelo identificou a solução ótima, consistente com a função objetivo e com todas as restrições impostas. A análise do conjunto de arestas demonstra, ainda, que o fluxo se distribuiu de maneira válida pela estrutura final, validando tanto a conservação quanto a utilização do fluxo como mecanismo de conectividade.

Em síntese, o trabalho realizado demonstra que a formulação por fluxo constitui não apenas uma forma exata e eficiente de resolver a AGM, mas também um arcabouço flexível, capaz de acomodar extensões e variações do problema original. A combinação entre rigor matemático, clareza estrutural e eficiência computacional torna essa abordagem particularmente adequada tanto para ambientes acadêmicos quanto para aplicações reais em áreas como redes de computadores, sistemas logísticos, planejamento de infraestrutura e engenharia de produção. Dessa forma, este estudo evidencia o potencial e a relevância da Programação Linear Inteira como ferramenta poderosa na solução de problemas de otimização sobre grafos.
